% GENERATED FILE. Edit canonical lessons, then run npm run generate:books.
% canonical-source-hash: fnv1a64:50e8ebb380080290
% canonical-lessons: MW-C101-yo, MW-C101-iso, MW-C101-uso, MW-C101-kiso, MW-C101-kun

\chapter{This, Like this, Like that, Which, Who}
\label{ch:mw-this-like-this-like-that-which-who}
\hlchaptermodality{101}

\begin{chapteropening}
\textbf{By the end of this chapter:} \emph{I can say this, like this, like that, which and who.}

Complete the last lesson of chapter 101: I can say this, like this, like that, which and who.
\end{chapteropening}

\section[\texorpdfstring{\mw{यो} (yo)}{यो (yo)}]{\mw{यो} (yo) — this}

\label{lesson:MW-C101-yo}



\begin{quote}
Before the new one: say the Marwadi for a fair, then the Marwadi for a weekly market.
\end{quote}

\subsection*{You'll want to know: \mw{यो}}
\textbf{\mw{यो}} — \emph{yo} — \textquotedblleft{}this\textquotedblright{}.

\begin{cousinweb}
CDIAL lists Hindi \emph{yah}, \textquotedblleft{}this\textquotedblright{}, under Sanskrit \emph{eṣa}, \textquotedblleft{}this\textquotedblright{} (CDIAL 2530). It gives no Marwari form, so the exact path to \textbf{\mw{यो}} (\emph{yo}) is not settled.
\end{cousinweb}

\begin{grammarlens}[title={What you've built}]
One more word for this chapter.
\end{grammarlens}

\subsection*{Your turn}
\begin{itemize}
\raggedright
  \item \emph{Say it:} \emph{yo}
  \item \emph{Say it:} \emph{yo}, once more
  \item \emph{Say it:} say \emph{yo}
  \item \emph{Recall:} say the Marwadi for a fair, then the Marwadi for a weekly market, then say \emph{yo} again
\end{itemize}

\begin{tcolorbox}[breakable,colback=teal!4,colframe=teal!35!black,title={Before you move on}]
What does \mw{यो} mean? (\textquotedblleft{}This\textquotedblright{}.) Say it once more.
\end{tcolorbox}
\section[\texorpdfstring{\mw{इसो} (iso)}{इसो (iso)}]{\mw{इसो} (iso) — like this, such}

\label{lesson:MW-C101-iso}



\begin{quote}
Before the new one: say the Marwadi for a weekly market, then the Marwadi for this.
\end{quote}

\subsection*{You'll want to know: \mw{इसो}}
\textbf{\mw{इसो}} — \emph{iso} — \textquotedblleft{}like this, such\textquotedblright{}.

\begin{cousinweb}
\textbf{\mw{इसो}} (\emph{iso}) is a tadbhava, a word inherited from Sanskrit and reshaped by sound change. It comes from \emph{īdṛśa}, \textquotedblleft{}such, like this\textquotedblright{}. CDIAL lists Old Marwari \emph{iso} under it (CDIAL 1611).
\end{cousinweb}

\begin{grammarlens}[title={What you've built}]
One more word for this chapter.
\end{grammarlens}

\subsection*{Your turn}
\begin{itemize}
\raggedright
  \item \emph{Say it:} \emph{iso}
  \item \emph{Say it:} \emph{iso}, once more
  \item \emph{Say it:} say \emph{iso}
  \item \emph{Recall:} say the Marwadi for a weekly market, then the Marwadi for this, then say \emph{iso} again
\end{itemize}

\begin{tcolorbox}[breakable,colback=teal!4,colframe=teal!35!black,title={Before you move on}]
What does \mw{इसो} mean? (\textquotedblleft{}Like this, such\textquotedblright{}.) Say it once more.
\end{tcolorbox}
\section[\texorpdfstring{\mw{उसो} (uso)}{उसो (uso)}]{\mw{उसो} (uso) — like that}

\label{lesson:MW-C101-uso}



\begin{quote}
Before the new one: say the Marwadi for this, then the Marwadi for like this.
\end{quote}

\subsection*{You'll want to know: \mw{उसो}}
\textbf{\mw{उसो}} — \emph{uso} — \textquotedblleft{}like that\textquotedblright{}.

\begin{cousinweb}
CDIAL has no entry for \textbf{\mw{उसो}} (\emph{uso}), so its history is not settled. It pairs with \emph{iso}, \textquotedblleft{}like this\textquotedblright{}, which CDIAL lists under Sanskrit \emph{īdṛśa}, \textquotedblleft{}such\textquotedblright{} (CDIAL 1611).
\end{cousinweb}

\begin{grammarlens}[title={What you've built}]
One more word for this chapter.
\end{grammarlens}

\subsection*{Your turn}
\begin{itemize}
\raggedright
  \item \emph{Say it:} \emph{uso}
  \item \emph{Say it:} \emph{uso}, once more
  \item \emph{Say it:} say \emph{uso}
  \item \emph{Recall:} say the Marwadi for this, then the Marwadi for like this, then say \emph{uso} again
\end{itemize}

\begin{tcolorbox}[breakable,colback=teal!4,colframe=teal!35!black,title={Before you move on}]
What does \mw{उसो} mean? (\textquotedblleft{}Like that\textquotedblright{}.) Say it once more.
\end{tcolorbox}
\section[\texorpdfstring{\mw{किसो} (kiso)}{किसो (kiso)}]{\mw{किसो} (kiso) — which, what kind of}

\label{lesson:MW-C101-kiso}



\begin{quote}
Before the new one: say the Marwadi for like this, then the Marwadi for like that.
\end{quote}

\subsection*{You'll want to know: \mw{किसो}}
\textbf{\mw{किसो}} — \emph{kiso} — \textquotedblleft{}which, what kind of\textquotedblright{}.

\begin{cousinweb}
\textbf{\mw{किसो}} (\emph{kiso}) is a tadbhava, a word inherited from Sanskrit and reshaped by sound change. It comes from \emph{kīdṛśa}, \textquotedblleft{}of what kind\textquotedblright{}. CDIAL lists Old Marwari \emph{kiso} under it (CDIAL 3197).
\end{cousinweb}

\begin{grammarlens}[title={What you've built}]
One more word for this chapter.
\end{grammarlens}

\subsection*{Your turn}
\begin{itemize}
\raggedright
  \item \emph{Say it:} \emph{kiso}
  \item \emph{Say it:} \emph{kiso}, once more
  \item \emph{Say it:} say \emph{kiso}
  \item \emph{Recall:} say the Marwadi for like this, then the Marwadi for like that, then say \emph{kiso} again
\end{itemize}

\begin{tcolorbox}[breakable,colback=teal!4,colframe=teal!35!black,title={Before you move on}]
What does \mw{किसो} mean? (\textquotedblleft{}Which, what kind of\textquotedblright{}.) Say it once more.
\end{tcolorbox}
\section[\texorpdfstring{\mw{कुण} (kuṇ)}{कुण (kuṇ)}]{\mw{कुण} (kuṇ) — who}

\label{lesson:MW-C101-kun}



\begin{quote}
Before the new one: say the Marwadi for like that, then the Marwadi for which.
\end{quote}

\subsection*{You'll want to know: \mw{कुण}}
\textbf{\mw{कुण}} — \emph{kuṇ} — \textquotedblleft{}who\textquotedblright{}.

\begin{cousinweb}
\textbf{\mw{कुण}} (\emph{kuṇ}) is a tadbhava, a word inherited from Sanskrit and reshaped by sound change. It comes from \emph{kaḥ punar}, \textquotedblleft{}who\textquotedblright{}. CDIAL lists Old Marwari \emph{kuṇa} under it (CDIAL 2575). It is built on \emph{kaḥ}, \textquotedblleft{}who\textquotedblright{}, and \emph{punar}, \textquotedblleft{}again\textquotedblright{}.
\end{cousinweb}

\begin{grammarlens}[title={What you've built}]
One more word for this chapter.
\end{grammarlens}

\subsection*{Your turn}
\begin{itemize}
\raggedright
  \item \emph{Say it:} \emph{kuṇ}
  \item \emph{Say it:} \emph{kuṇ}, once more
  \item \emph{Say it:} say \emph{kuṇ}
  \item \emph{Recall:} say the Marwadi for like that, then the Marwadi for which, then say \emph{kuṇ} again
\end{itemize}

\begin{tcolorbox}[breakable,colback=teal!4,colframe=teal!35!black,title={Before you move on}]
What does \mw{कुण} mean? (\textquotedblleft{}Who\textquotedblright{}.) Say it once more.
\end{tcolorbox}
